\section{Technical Feasibility}
The proposed machine uses proven mechanical subsystems (indexing mechanism, feed guides, and a binding station) that can be manufactured with conventional workshop processes. Critical components such as cams, shafts, bearings and fasteners are standard items, minimizing custom machining.

\begin{table}[h]
\centering
\caption{Representative Hardware Specifications}
\begin{tabular}{|l|l|}
\hline
Component & Typical Specification \\
\hline
Frame & Mild steel sheet / welded frame \\
Indexing mechanism & Geneva or cam-driven with hardened pins \\
Bearings & Standard ball or sleeve bearings (sizes as per shaft) \\
Feed guides & Low-friction polymer or brass inserts \\
Drive & Hand-crank or small geared motor (10–50 W) \\
Binding mechanism & Simple clamp or adhesive dispenser \\
Fasteners & Standard bolts, nuts, washers \\
\hline
\end{tabular}
\label{tab:hardware-specs}
\end{table}

\section{Economic Feasibility}
Initial material and fabrication costs are estimated to be low-to-moderate due to the use of off-the-shelf components and simple manufacturing methods. The design targets affordable parts and minimizes electronic control to reduce both capital cost and maintenance burden.

\section{Operational Feasibility}
The machine is designed for ease of use and maintenance. Operators require minimal training to load sticks, clear jams, and perform routine lubrication. Spare parts are inexpensive and widely available.

\section{Schedule and Deliverables}
A prototype build, bench testing, and iterative refinement are planned over an expected timeline of 8–12 weeks: detailed design (2–3 weeks), fabrication (3–4 weeks), testing and validation (3–5 weeks).

